\chapter{Annexes}
\label{chap:annexes}

\section{Raccourcis clavier}

\begin{tabularx}{\linewidth}{l X}
    \toprule
    \textbf{Raccourci} & \textbf{Action} \\
    \midrule
    \touche{Ctrl}+\touche{O}            & Ouvrir profil courant \\
    \touche{Ctrl}+\touche{Maj}+\touche{O} & Ouvrir profil de référence \\
    \touche{Ctrl}+\touche{S}            & Sauvegarder profil courant \\
    \touche{Ctrl}+\touche{Q}            & Quitter \\
    \touche{Ctrl}+\touche{Z}            & Annuler (réservé, non actif) \\
    \touche{Ctrl}+\touche{Y}            & Refaire (réservé, non actif) \\
    \touche{Ctrl}+\touche{B}            & Convertir en spline \\
    \touche{Ctrl}+\touche{0}            & Zoom adapté (canvas profils) \\
    \bottomrule
\end{tabularx}

\section{Glossaire}

\begin{description}
    \item[BA / BF] Bord d'attaque ($x=0$) / Bord de fuite ($x=c$).
    \item[Cambrure] Hauteur maximale de la ligne moyenne du profil,
        exprimée en pourcentage de la corde.
    \item[$C_D$] Coefficient de traînée. Force de traînée
        normalisée par la pression dynamique et la corde.
    \item[$C_f$] Coefficient de frottement local pariétal.
    \item[$C_L$] Coefficient de portance.
    \item[$C_M$] Coefficient de moment de tangage, généralement
        référencé au quart de corde.
    \item[$C_p$] Coefficient de pression : $(p - p_\infty) /
        (\frac{1}{2}\rho V_\infty^2)$.
    \item[Continuité C\textsuperscript{0}, C\textsuperscript{1},
        C\textsuperscript{2}] Au point de raccord de deux courbes :
        identité de position (C\textsuperscript{0}), de tangente
        (C\textsuperscript{1}), de courbure (C\textsuperscript{2}).
    \item[Corde] Distance entre le BA et le BF, sert de longueur
        de référence.
    \item[Courbe de Bézier] Courbe paramétrique définie par un
        polygone de contrôle. Le degré est égal au nombre de
        points de contrôle moins un.
    \item[$\delta^*$] Épaisseur de déplacement de la couche limite.
    \item[De Casteljau] Algorithme récursif d'évaluation et de
        subdivision d'une courbe de Bézier. La subdivision est
        \emph{exacte} (pas d'approximation).
    \item[Décrochage] Perte de portance brutale au-delà d'une
        incidence critique, due à un décollement de la couche
        limite.
    \item[Extrados / Intrados] Surface supérieure / inférieure du
        profil.
    \item[Finesse] Rapport $C_L/C_D$. Plus la finesse est élevée,
        plus le rendement aérodynamique est bon.
    \item[$H$] Facteur de forme de la couche limite,
        $H = \delta^*/\theta$.
    \item[NCRIT] Critère de transition par méthode $e^N$ ;
        amplitude maximale acceptée des perturbations avant
        déclenchement de la turbulence.
    \item[Polaire] Tracé d'une grandeur aérodynamique en
        fonction d'une autre, à un Reynolds donné. Plusieurs
        polaires forment une famille indexée par le Reynolds.
    \item[Reynolds] Nombre adimensionnel
        $\mathrm{Re} = \rho V c / \mu$, caractérisant le régime
        d'écoulement (laminaire, transitionnel, turbulent).
    \item[Selig] Convention de stockage des points d'un profil :
        BF $\to$ extrados $\to$ BA $\to$ intrados $\to$ BF.
    \item[Spline] Courbe par morceaux composée de plusieurs
        segments de Bézier raccordés selon un niveau de continuité
        donné.
    \item[$\theta$] Épaisseur de quantité de mouvement de la
        couche limite.
    \item[Trip / Turbulator] Élément (rugosité, fil tendu) forçant
        la transition laminaire-turbulent à une position donnée.
    \item[$U_e$] Vitesse au bord externe de la couche limite.
    \item[XFoil] Solveur aérodynamique 2D développé par Mark Drela
        (MIT) couplant méthode des panneaux et équations intégrales
        de couche limite.
    \item[XTR] Position de transition de couche limite ($x_{tr}/c$).
\end{description}

\section{Dépannage (FAQ)}

\subsection*{La GUI ne se lance pas}

Vérifiez que :

\begin{enumerate}
    \item l'environnement virtuel est activé (\chemin{env\_py3}) ;
    \item les dépendances sont installées (\code{pip install -r
        requirements.txt}) ;
    \item Python est en version 3.10 ou supérieure
        (\code{python --version}).
\end{enumerate}

\subsection*{XFoil ne s'exécute pas}

\begin{itemize}
    \item Vérifiez la présence de \chemin{externaltools/xfoil/xfoil.exe}.
    \item Sur Windows, certains antivirus peuvent bloquer
        l'exécution. Ajoutez une exception pour le dossier
        \chemin{externaltools}.
    \item Les commandes envoyées à XFoil sont visibles dans la
        console (en mode développement) ; consultez la sortie
        pour diagnostiquer une éventuelle erreur de syntaxe.
\end{itemize}

\subsection*{XFoil ne converge pas pour certains angles}

\begin{itemize}
    \item Désactivez \menu{ASEQ} pour passer en mode ALFA
        bidirectionnel.
    \item Réduisez le pas en $\alpha$ (par exemple de 0{,}5\textdegree{} à
        0{,}25\textdegree{}) pour faciliter le suivi de continuation.
    \item Augmentez \menu{Timeout} si le calcul est long.
    \item Vérifiez que le profil ne présente pas de singularités
        géométriques (points dupliqués, croisement extrados/intrados,
        BA dégénéré).
\end{itemize}

\subsection*{La courbure ne s'affiche pas}

\begin{itemize}
    \item Convertissez le profil en mode spline
        (\menu{Édition $\rightarrow$ Convertir en Spline}).
    \item Augmentez l'échelle des porcupines (clic droit sur le
        canvas $\rightarrow$ \menu{Échelle porcupines\dots}) si
        elles sont trop courtes pour être visibles.
\end{itemize}

\subsection*{Je perds mes splines en sauvegardant}

Le format \chemin{.dat} (Selig/Lednicer) ne préserve pas la
représentation spline : seuls les points discrets sont
sauvegardés. Pour conserver la spline, utilisez le format
\chemin{.bspl}.

\section{Pour aller plus loin}

\begin{itemize}
    \item \textbf{Documentation technique} : voir le fichier
        \chemin{CLAUDE.md} à la racine du dépôt et les
        docstrings du code source (format Sphinx reST).
    \item \textbf{Tests unitaires} : voir le dossier
        \chemin{tests/} qui contient plus de 240 tests
        (Bezier, Profil, BezierSpline, ProfilSpline,
        Simulation).
    \item \textbf{Article de référence sur XFoil} :
        Drela, M. \emph{XFOIL: An Analysis and Design System for
        Low Reynolds Number Airfoils}, in \emph{Low Reynolds
        Number Aerodynamics}, Springer Lecture Notes in Engineering,
        Vol.~54, 1989.
    \item \textbf{Bibliothèque de profils} :
        \chemin{https://airfoiltools.com/} (large catalogue
        au format \chemin{.dat}).
    \item \textbf{Site Nervures} :
        \chemin{https://www.nervures.com/}.
\end{itemize}

\bigskip
\hrule
\bigskip

\begin{center}
    \itshape
    Fin du manuel utilisateur.\\
    Pour signaler une coquille, suggérer une amélioration ou
    poser une question, ouvrez une \emph{issue} sur
    \chemin{https://github.com/BenoitGFree/AifoilTools}.
\end{center}
